% ===================================================================================
% Section 5 (chunk C): the Riemann--Hilbert problem and the common frame.
% Labels defined here: sec:frame, sec:frame-*, lem:interp, lem:bandweight, lem:gfunction,
% lem:signs, lem:weightdomain, lem:pottersector, lem:jumpsV, tab:geometry, eq:*, fig:geometry.
% External labels used: prop:equilibrium (Prop. 2.2), lem:potter (Lemma 2.1), tool:T1..T4 (Sect. 2.5).
% ===================================================================================
\providecommand{\eps}{\varepsilon}
\providecommand{\Log}{\operatorname{Log}}
\providecommand{\re}{\operatorname{Re}}
\providecommand{\im}{\operatorname{Im}}
\providecommand{\Res}{\operatorname*{Res}}
\providecommand{\thlip}{\vartheta_{\mathrm{lip}}}
\providecommand{\thseam}{\vartheta_{\mathrm{seam}}}
\providecommand{\thw}{\vartheta_{\mathrm{w}}}
\providecommand{\Ob}{\Omega_{\mathrm{b}}}
\providecommand{\Osat}{\Omega^{\triangle}}
\providecommand{\Oinf}{\Omega_{\infty}}
\providecommand{\Sin}{S_{\mathrm{in}}}
\providecommand{\Sout}{S_{\mathrm{out}}}
\providecommand{\Zn}{\mathcal{Z}_n}
\providecommand{\Nn}{\mathcal{N}}
\providecommand{\tvphi}{\tilde\varphi}
\providecommand{\Fi}[1]{\textup{(F#1)}}

\section{The Riemann--Hilbert problem and the common frame}\label{sec:frame}

This section and the next prove the part of Theorem~\ref{thm:engine} that does not concern the hard edge. We
encode the discrete measure by a meromorphic interpolation problem of Fokas--Its--Kitaev type, rescale it by $n^p$,
and normalise it with the $g$-function of the equilibrium measure $\mu^*$ of Paper~1. After rescaling, the nodes
fill $(0,n^p]$. On the band $(0,4c)$ the orthogonal polynomials oscillate; on the saturated region $(4c,n^p)$ the
nodes are as dense as the equilibrium measure allows, and the problem is close to a pure interpolation problem;
near $0$ the band meets infinitely many nodes, and $0$ is a hard edge. The frame is that of Wu, Lin, Xu and Zhao
for the Tricomi--Carlitz polynomials (\ref{tool:T1}), written in the one-sided variable $x$.

Almost every step is the same for all $p>0$. The only geometry-dependent ingredient is a scalar function $F$, the
\emph{outer scalar}, which removes the Gamma-function growth of the node masses away from the hard edge. We list
in~\S\ref{sec:frame-outer} the five properties of $F$ that the analysis uses, and we construct $F$ in
\S\ref{sec:outer} separately for $p>1$, $p=1$ and $0<p<1$. Table~\ref{tab:geometry} at the end of this section
collects everything that depends on the geometry.

\subsection{Standing assumptions}\label{sec:frame-standing}
Throughout \S\S\ref{sec:frame}--\ref{sec:outer} and Appendix~\ref{app:apriori} we fix $p>0$ and $1<s<1+2p$, and we
put
\begin{equation}\label{eq:alpha}
  \alpha:=\frac{s-1-p}{p}\in(-1,1).
\end{equation}
We fix $\ell_1(t)=e^{\Lambda(\log t)}$ satisfying $(\mathrm A_{\vartheta_p})$, with the strip half-width
$\vartheta_p=\pi/p$ for $p>1$ and $\vartheta_p=\pi/2$ for $0<p\le1$, and with the constant $c_0>0$ of that
hypothesis. We write $\sigma:=\sigma_{p,s,\ell_1}=\sum_{k\ge1}w_k\delta_{x_k}$, with $x_k=k^{-p}$ and
$w_k=k^{-s}\ell_1(k)$, and $c=c(p)$, $\gamma_p=\log c+p$. The degree $n$ is large, $n\ge n_0$, where $n_0$ may
increase from line to line. Constants $C,c',\dots$ (never the letter $c$, which is $c(p)$) depend only on $p$, $s$,
$\ell_1$ and the contour parameters fixed in~\S\ref{sec:frame-contour}; they do not depend on $n$ unless they carry
the subscript $n$. We put
\begin{equation}\label{eq:thetaw}
  \thw:=\begin{cases}7\pi/16,&p\ge1,\\ 3\pi p/8,&0<p<1,\end{cases}\qquad\qquad
  \thlip:=\begin{cases}\pi/3,&p\ge1,\\ \pi p/4,&0<p<1,\end{cases}
\end{equation}
so that $0<\thlip<\thw$ and $\thw/p<\vartheta_p$ in all three cases ($\thw/p=7\pi/(16p)<\pi/p$ for $p>1$, $\thw=7\pi/16<\pi/2$ for $p=1$, and
$\thw/p=3\pi/8<\pi/2$ for $p<1$); moreover $\frac54\thlip<\thw$, and for $p<1$, $\thw<\pi p/2$. The angle $\thw$ bounds the arguments at which the weight is ever evaluated
(Lemma~\ref{lem:weightdomain}); $\thlip$ is the angle at which the lens leaves the hard edge.

\subsection{The interpolation problem}\label{sec:frame-interp}
Let $P_n$ be the monic orthogonal polynomials of $\sigma$ and $h_n=\int P_n^2\diff\sigma$. For a $2\times2$ matrix
function $X$ with an expansion $X(z)=I+X_1/z+X_2/z^2+O(z^{-3})$ at $\infty$ we write
\begin{equation}\label{eq:functional-a}
  \mathsf a(X):=\frac{(X_2)_{12}}{(X_1)_{12}}-(X_1)_{22}\qquad\text{(when $(X_1)_{12}\ne0$).}
\end{equation}

\begin{lemma}[The interpolation problem]\label{lem:interp}
Let $Y$ be the matrix \eqref{eq:Y-solution} of \ref{tool:T1}(a), built from the monic orthogonal polynomials $P_n$,
$P_{n-1}$ of $\sigma$ and $h_{n-1}=\int P_{n-1}^2\diff\sigma$. Then $\det Y\equiv1$. The rescaled matrix
\begin{equation}\label{eq:U}
  U(z):=n^{pn\sigma_3}\,Y(zn^{-p})
\end{equation}
is analytic off $\{X_k\}_{k\ge1}\cup\{0\}$, where
\begin{equation}\label{eq:Xk}
  X_k:=n^px_k=(n/k)^p\quad(k\ge1),\qquad X_1=n^p\ \text{(the top node)};
\end{equation}
it has simple poles at the $X_k$ in its second column only, with $\Res_{z=X_k}U_{j2}=n^pw_k\,U_{j1}(X_k)$ ($j=1,2$);
and $U(z)z^{-n\sigma_3}=I+U_1/z+U_2/z^2+\dots$ for $|z|>n^p$, with
\begin{equation}\label{eq:extractU}
  (U_1)_{12}(U_1)_{21}=n^{2p}b_n^2,\qquad \mathsf a(Uz^{-n\sigma_3})=n^pa_n .
\end{equation}
\end{lemma}
% src: RH_px.md §1, §5 "Extraction"; WLXZ Thm 2.3 (T1); review REVIEW_RHpx "Extraction indices"

\begin{proof}
For the determinant, write $h_{n-1}\det Y$ as
\begin{align*}
  \int\frac{P_n(z)P_{n-1}(x)-P_{n-1}(z)P_n(x)}{z-x}\diff\sigma(x)
  &=\int\frac{P_n(z)-P_n(x)}{z-x}P_{n-1}(x)\diff\sigma(x)\\
  &\quad-\int\frac{P_{n-1}(z)-P_{n-1}(x)}{z-x}P_n(x)\diff\sigma(x).
\end{align*}
The first difference quotient is a polynomial in $x$ of degree $n-1$ with leading coefficient $1$, so the first
integral is $h_{n-1}$; the second has degree $n-2$ in $x$, so the second integral vanishes. Hence $\det Y\equiv1$.
The statements for $U$ follow from those of \ref{tool:T1}(a) by the substitution $z=n^p\zeta$: the residue picks up
the factor $n^p$ from $\diff z=n^p\diff\zeta$, and $U_j=n^{jp}\,n^{pn\sigma_3}Y_jn^{-pn\sigma_3}$ ($j=1,2$). The
conjugation multiplies the $(1,2)$ entries by $n^{2pn}$ and the $(2,1)$ entries by $n^{-2pn}$; this cancels in the
product and in the ratio in~\eqref{eq:functional-a}, and \eqref{eq:Y-extract} gives~\eqref{eq:extractU}.
\end{proof}

We never use that $Y$ is the \emph{unique} solution of an interpolation problem: every object below is built from
$Y$, and only its explicit properties are used. In particular no existence theorem for a Riemann--Hilbert problem
is needed.

\subsection{The index function and the band weight}\label{sec:frame-weight}
The node set $\{X_k\}$ is the set of points where the \emph{index function}
\begin{equation}\label{eq:t}
  t(z):=nz^{-1/p}\qquad(z\in\C\setminus(-\infty,0],\ \text{principal power})
\end{equation}
takes positive integer values: $t(X_k)=k$. We put
\begin{equation}\label{eq:theta}
  \theta(z):=\pi t(z),\qquad \Pi(z):=\frac{\sin\theta(z)}{\theta'(z)},\qquad \theta'(z)=-\frac{\pi}{p}\,\frac{t(z)}{z},
\end{equation}
so that $\Pi$ has simple zeros at the nodes. On $\C\setminus(-\infty,0]$ we have $\arg t=-\arg z/p$, hence
\begin{equation}\label{eq:imtheta}
  \pm\im\theta(z)<0\quad\text{for }0<\pm\arg z<\min(\pi,\pi p),\qquad
  |e^{\mp2i\theta(z)}|=\exp\bigl(-2\pi n|z|^{-1/p}\sin(|\arg z|/p)\bigr).
\end{equation}
The weight enters through the log variable $\xi:=\log t=\log n-(\Log z)/p$. Let
\begin{equation}\label{eq:Zn}
  \Zn:=\{z:\ 0<|z|<e^{pc_0}n^p,\ |\arg z|<\thw\}.
\end{equation}
For $z\in\Zn$ we have $|\im\xi|<\thw/p<\vartheta_p$ and $\re\xi>-c_0$, so $\xi\in S_{\vartheta_p}$ and the
\emph{residue weight} and the \emph{slowly varying factor}
\begin{equation}\label{eq:w-rn}
  w(z):=n^p\,t(z)^{-s}\,\ell_1(t(z)),\qquad r_n(z):=\frac{\ell_1(t(z))}{\ell_1(n)}=e^{\Lambda(\xi)-\Lambda(\log n)}
\end{equation}
are analytic on $\Zn$. By Lemma~\ref{lem:interp}, $\Res_{X_k}U_{j2}=w(X_k)U_{j1}(X_k)$. The combination that appears
in all jumps is
\begin{equation}\label{eq:W}
  W(z):=2i\theta'(z)w(z)\qquad(z\in\Zn).
\end{equation}

\begin{lemma}[Band weight]\label{lem:bandweight}
On $\Zn$,
\begin{equation}\label{eq:Wformula}
  W(z)=C_n'\,z^{\alpha}\,r_n(z),\qquad C_n':=-\frac{2\pi i}{p}\,n^{1+p-s}\ell_1(n),
\end{equation}
with the principal power. In particular $|W(x)|=\frac{2\pi}{p}n^{1+p-s}\ell_1(n)\,x^\alpha r_n(x)$ and
$\arg W(x)=-\pi/2$ for $x\in(0,e^{pc_0}n^p)$, where $r_n(x)>0$. The weight has a hard edge at $0$ with Bessel order
$\alpha$, and $\alpha>-1$ if and only if $s>1$.
\end{lemma}
% src: RH_px Lemma 1.2; ROUTE2 §2 table row "Lemma 1.2"; xframe §2 row L1.2

\begin{proof}
By~\eqref{eq:theta}, $2i\theta'w=-\frac{2\pi i}{p}\frac{t}{z}n^pt^{-s}\ell_1(t)$, and $t^{1-s}=n^{1-s}z^{(s-1)/p}$.
The exponent of $z$ is $(s-1)/p-1=\alpha$. Finally $\ell_1(t)=\ell_1(n)r_n$, and $r_n>0$ on $(0,\infty)$ because
$\Lambda$ is real on the real axis.
\end{proof}

\subsection{The $g$-function}\label{sec:frame-g}
We use the equilibrium measure of Paper~1 exactly as quoted in Proposition~\ref{prop:equilibrium}: $\mu^*=\int_0^1
\omega_{\beta_\tau}\diff\tau$ with $\beta_\tau:=4c\tau^{-p}$ and $\omega_\beta$ the arcsine law on $[0,\beta]$; its
density $\psi^*$ satisfies $\psi^*=\lambda_p$ on $[4c,\infty)$ and $0<\psi^*<\lambda_p$ on $(0,4c)$, where
$\lambda_p(x)=\frac1px^{-1-1/p}$ is the node density; and $U^*(x)=\int\log|x-y|\diff\mu^*(y)$ equals $\gamma_p$ on
$(0,4c]$. Note $\int_x^\infty\lambda_p=x^{-1/p}=t(x)/n$. Put
\begin{equation}\label{eq:g}
  g(z):=\int\log(z-y)\diff\mu^*(y)\quad(z\in\C\setminus\R),\qquad
  \varphi:=\gamma_p-g,\qquad \tvphi:=\varphi\pm\frac{i\theta}{n}\ \text{ in }\C_\pm,
\end{equation}
with principal logarithms. Nothing here depends on $\ell_1$: the external field of the weight is
$\frac1n\log w=O(\frac{\log n}{n})$ on the contour (Lemma~\ref{lem:pottersector}) and is not put into $g$.

\begin{lemma}[$g$-function]\label{lem:gfunction}
\begin{enumerate}
\item[(i)] $g$ has continuous boundary values $g_\pm$ on $\R\setminus\{0\}$, $g_+-g_-=2\pi i$ on $(-\infty,0)$ and
  $g_+-g_-=2\pi i\mu^*((x,\infty))$ on $(0,\infty)$. In particular $\varphi_++\varphi_-=0$ on $(0,4c)$, and $\tvphi$
  has no jump on $(4c,\infty)$.
\item[(ii)] $g(z)-\Log z\to0$ as $z\to\infty$, uniformly in $\C\setminus\R$.
\item[(iii)] $U^*$ is strictly increasing on $[4c,\infty)$, and $U^*(x)=\log x+o(1)$ as $x\to\infty$. Hence
  $\re\tvphi(x)=\gamma_p-U^*(x)<0$ for $x>4c$, and $\re\tvphi(x)\le-\frac12\log x$ for $x\ge x_0$.
\end{enumerate}
\end{lemma}
% src: RH_px Lemma 2.1 and §4.3; review REVIEW_RHpx (strictness); REVIEW_Plan54 F1 (rate at p = 1)

\begin{proof}
(i) For $y>0$, $\log(z-y)$ jumps by $2\pi i$ across $(-\infty,y)$. Integrating against $\mu^*$ gives the jumps;
continuity of the boundary values follows from the continuity of $\psi^*$ on $(0,\infty)$. On $(0,4c)$,
$g_++g_-=2U^*=2\gamma_p$, so $\varphi_++\varphi_-=0$. On $(4c,\infty)$, $\mu^*((x,\infty))=\int_x^\infty\lambda_p=
t(x)/n$, so $g_+-g_-=2i\theta/n$, and $\tvphi_+-\tvphi_-=-(g_+-g_-)+2i\theta/n=0$.

(ii) For $z\in\C_+$ and $y>0$ we have $\Log(z-y)=\Log z+\Log(1-y/z)$, because $\arg(z-y)-\arg z\in(0,\pi)$; the
case $z\in\C_-$ is the mirror image. So $g(z)-\Log z=\int\Log(1-y/z)\diff\mu^*(y)$. On $\{y\le|z|/2\}$ we use
$|\Log(1-y/z)|\le2y/|z|$ and $\int_{y\le R}y\diff\mu^*\le1+\int_1^Ry\lambda_p(y)\diff y=o(R)$ (as $\psi^*\le\lambda_p$). On
$\{y>|z|/2\}$ we use $\psi^*\le\lambda_p\le C|z|^{-1-1/p}$ there and
$|\Log(1-y/z)|\le\pi+\log(1+u)+|\log|1-u||$ with $u=y/|z|$, which follows from
$|1-u|\le|1-y/z|\le1+u$. The resulting integral is $O(|z|^{-1/p})$. Both parts tend to $0$ uniformly.

(iii) For $\beta>0$ put $G_\beta(z):=2\log\bigl((\sqrt z+\sqrt{z-\beta})/2\bigr)$. It is analytic off
$(-\infty,\beta]$, equals $\log z+o(1)$ at $\infty$, and $G_\beta'(z)=1/\sqrt{z(z-\beta)}=\int\diff\omega_\beta(y)/(z-y)$,
the Stieltjes transform of the arcsine law. Hence $G_\beta(z)=\int\log(z-y)\diff\omega_\beta(y)$. Taking real parts,
$u_\beta(x):=\int\log|x-y|\diff\omega_\beta(y)$ equals $\log(\beta/4)$ on $[0,\beta]$ (there
$|\sqrt x+i\sqrt{\beta-x}|^2=\beta$), and it is strictly increasing on $[\beta,\infty)$. So
$U^*=\int_0^1u_{\beta_\tau}\diff\tau$ is non-decreasing on $(0,\infty)$. For $x>4c$ the set of $\tau$ with
$\beta_\tau<x$ is the interval $((4c/x)^{1/p},1]$, of positive length, on which $u_{\beta_\tau}$ is strictly
increasing at $x$. So $U^*$ is strictly increasing on $[4c,\infty)$. For the asymptotics, $U^*(x)-\log
x=\int\log|1-y/x|\diff\mu^*(y)$, and the argument of~(ii) on the real axis shows that this tends to $0$. Finally,
$\re\tvphi=\re\varphi=\gamma_p-U^*$ on $(4c,\infty)$, because $\theta$ is real there; $U^*(4c)=\gamma_p$, which gives
the sign, and $U^*(x)\ge\frac12\log x+\gamma_p$ for large $x$.
\end{proof}

At $p=1$ the first moment of $\mu^*$ diverges and $g(z)-\Log z=O(|z|^{-1}\log|z|)$; for $p>1$ it is
$O(|z|^{-1/p})$ and for $p<1$ it is $O(|z|^{-1})$. Only~(ii) is used.

\subsection{Signs and the local variables at the two edges}\label{sec:frame-signs}
The lens opening needs $\re\varphi<0$ off the band, and the triangular factors that appear later need
$\re\tvphi>0$ just above and below the band. Near the edges we need the local variables of the Bessel and Airy
models.

\begin{lemma}[Signs and local variables]\label{lem:signs}
\begin{enumerate}
\item[(i)] $\re\varphi(z)<0$ for every $z\in\C\setminus[0,4c]$ (for $z$ on $(4c,\infty)$ and $(-\infty,0)$ this
  refers to the common real part of the two boundary values), and $\re\varphi_\pm=0$ on $(0,4c)$.
\item[(ii)] For every compact $K\subset(0,4c)$ there is $y_K>0$ such that $\re\tvphi(x\pm iy)>0$ for $x\in K$ and
  $0<y\le y_K$. Moreover $\partial_y\re\tvphi(x+iy)|_{y=0^+}=\pi(\lambda_p-\psi^*)(x)>0$ on $(0,4c)$.
\item[(iii)] (Hard edge.) There are $\delta_0>0$ and a function $f$, analytic and univalent on $U(0,2\delta_0)$,
  real on the real axis, with $f(0)=0$ and $f'(0)<0$, such that $(\varphi\pm i\pi)^2=4f$ on $U(0,2\delta_0)\cap\C_\pm$.
  More precisely $\varphi=\mp i\pi-2f^{1/2}$ in $\C_\pm$, with the principal root. In particular $f$ does not
  depend on $n$, $|\re\varphi(z)|\le C|z|^{1/2}$ on $U(0,\delta_0)$, and $\re\tvphi(z)\ge\pi|z|^{-1/p}\sin(|\arg
  z|/p)-C|z|^{1/2}$ for $z\in U(0,\delta_0)$, $0<|\arg z|\le\thw$.
\item[(iv)] (Soft edge.) There are $\eps_0>0$ and a function $f_r$, analytic and univalent on $U(4c,2\eps_0)$, real
  on the real axis, with $f_r(4c)=0$ and $f_r'(4c)>0$, such that $\tvphi^2=\frac49f_r^3$ on $U(4c,2\eps_0)$, and
  $\tvphi=-\frac23f_r^{3/2}$ (principal branch, cut along $f_r^{-1}((-\infty,0])$). In particular
  $\re\tvphi>0$ on $\{z\in U(4c,\eps_0):\ \pi/3<|\arg(z-4c)|<\pi\}$ and $\re\tvphi<0$ for $|\arg(z-4c)|<\pi/3$.
\end{enumerate}
\end{lemma}
% src: RH_px Lemma 2.2, Lemma 4.1, §4.2; review REVIEW_RHpx ("Near 0, Re phi < 0 everywhere off the band")

\begin{proof}
(i) We have $\re\varphi(z)=\gamma_p-\int\log|z-y|\diff\mu^*(y)$. If $\re z>0$ and $\im z\ne0$, then
$|z-y|>|\re z-y|$ for all $y$, so $\re g(z)>U^*(\re z)\ge\gamma_p$ (Proposition~\ref{prop:equilibrium} and
Lemma~\ref{lem:gfunction}(iii)). On $(4c,\infty)$, $\re g_\pm=U^*>\gamma_p$ by Lemma~\ref{lem:gfunction}(iii). If
$\re z\le0$ and $z\ne0$, choose $0<\eps\le\min(|z|/2,4c)$; then $|z-y|\ge y>y-\eps$ for $y\ge\eps$, and $|z-y|\ge|z|>\eps\ge
|\eps-y|$ for $0<y<\eps$, so $\re g(z)>U^*(\eps)=\gamma_p$. On $(0,4c)$, $\re g_\pm=U^*=\gamma_p$.

(ii) In $\C_+$, $\tvphi=\varphi+i\theta/n=\gamma_p-g+i\pi z^{-1/p}$. Since $\psi^*$ is real-analytic on $(0,4c)$ (it is
$\lambda_p$ times a regularised incomplete beta function of $x/4c$), $g|_{\C_+}$ continues analytically across
$(0,4c)$, and so does $\tvphi$. By the Cauchy--Riemann equations, $\partial_y\re g(x+i0)=-\im g'(x+i0)=\pi\psi^*(x)$
and $\partial_y\re(i\pi z^{-1/p})|_{z=x}=-\pi\,\frac{\diff}{\diff x}x^{-1/p}=\pi\lambda_p(x)$. Since $\re\tvphi=0$ on
$(0,4c)$, a Taylor expansion with a uniform remainder on $K$ gives the claim. $\C_-$ is the mirror image.

(iii) On $(0,4c)$, $\varphi_+(x)=\gamma_p-U^*(x)-i\pi\mu^*((x,\infty))=-i\pi+i\pi\mu^*((0,x])$. By
Proposition~\ref{prop:equilibrium}, $\psi^*(x)=\lambda_p(x)\,\mathcal B_{x/4c}(\frac1p+\frac12,\frac12)$ near $0$,
and $\mathcal B_y(a,\frac12)=y^aA_0(y)$ with $A_0$ analytic near $0$ and $A_0(0)>0$. Since
$\lambda_p(x)(x/4c)^{1/p+1/2}=Cx^{-1/2}$, we get $\mu^*((0,x])=x^{1/2}A(x)$ with $A$ analytic near $0$ and $A(0)>0$.
Put $f(z):=-\frac{\pi^2}{4}zA(z)^2$. Then $f$ is analytic and real on the real axis, $f(0)=0$, $f'(0)=-\pi^2A(0)^2/4<0$,
so $f$ is univalent on a disc $U(0,2\delta_0)$. The functions $(\varphi+i\pi)^2$ (on $\C_+$) and $4f$ are analytic in
$U(0,2\delta_0)\cap\C_+$ and have the same continuous boundary values $-\pi^2xA(x)^2$ on $(0,2\delta_0)$; by the
reflection principle they coincide. The same holds in $\C_-$ with $\varphi-i\pi$. On the band, $f(x+i0)$ lies on the
negative axis approached from below, and the principal root gives $f^{1/2}=-\frac{i\pi}2x^{1/2}A(x)$, so
$\varphi_++i\pi=-2f^{1/2}$; by continuity $\varphi=-i\pi-2f^{1/2}$ in $U(0,2\delta_0)\cap\C_+$. The remaining claims
follow from $|f(z)|\le C|z|$, from $\tvphi=\varphi\pm i\pi z^{-1/p}$ and from
$\re(\pm i\pi z^{-1/p})=\pi|z|^{-1/p}\sin(|\arg z|/p)$.

(iv) On $(4c-2\eps_0,4c)$ we have $\tvphi_++\tvphi_-=\varphi_++\varphi_-=0$ and, using
Lemma~\ref{lem:gfunction}(i), $\tvphi_+-\tvphi_-=2\pi i\int_x^\infty(\lambda_p-\psi^*)=2\pi i\int_x^{4c}(\lambda_p-\psi^*)$.
Near $y=1$, $1-\mathcal B_y(a,\frac12)=(1-y)^{1/2}A_1(y)$ with $A_1$ analytic and $A_1(1)>0$, so
$\tvphi_+(x)=i\pi(4c-x)^{3/2}A_2(x)$ with $A_2$ analytic near $4c$ and $A_2(4c)>0$. As in~(iii), $\tvphi^2$ and
$h(z):=\pi^2(z-4c)^3A_2(z)^2$ are analytic in $U(4c,2\eps_0)\setminus[4c-2\eps_0,4c+2\eps_0]$ with equal boundary
values on $(4c-2\eps_0,4c)$, and $\tvphi$ has no jump on $(4c,4c+2\eps_0)$ (Lemma~\ref{lem:gfunction}(i)). So
$\tvphi^2=h$. Put $f_r:=(9h/4)^{1/3}=(\frac{3\pi}{2})^{2/3}(z-4c)A_2(z)^{2/3}$. On $(4c,\infty)$, $\tvphi$ is real and
negative (Lemma~\ref{lem:gfunction}(iii)), which fixes the branch $\tvphi=-\frac23f_r^{3/2}$. The sign statements
follow from $\arg f_r^{3/2}=\frac32\arg(z-4c)+o(1)$.
\end{proof}

\subsection{The contour}\label{sec:frame-contour}
We fix, once and for all and independently of $n$, numbers $\delta,\eps,h_0,\delta'$ with
\[
  0<\eps<h_0<\eps_0,\qquad 0<\delta<\min(\delta_0,\tfrac14,\tfrac c2),\qquad h_0\le\tfrac12\delta\tan\thw,\qquad 0<\delta'<1,
\]
and the following arcs (Figure~\ref{fig:geometry}).
\begin{itemize}
\item The \emph{band} $[0,4c]$, the \emph{saturated segment} $(4c,n^p+\delta')$ and the \emph{top segment}
  $\Sigma_{\mathrm{top}}:=\{n^p+\delta'+iy:\ 0<|y|\le h_0\}$. The top node $X_1=n^p$ lies inside, and no node lies on
  $[n^p+\delta',\infty)$.
\item The \emph{lips} $\Sigma_{\mathrm{lip}}^\pm$, $C^2$ arcs in $\C_\pm$ from $0$ to $4c\pm ih_0$. In $U(0,\delta)$
  they are the preimages under $\zeta=n^2f(z)$ of the rays $\arg\zeta=\mp(\pi-\thlip)$; since $f'(0)<0$, they leave
  $0$ at the angles $\pm\thlip$, and they do not depend on $n$. Outside $U(0,\delta)$
  they run in $\C_\pm\cap\{|\arg z|<\thw\}$, inside the domain where $\re\tvphi>0$ (such arcs exist by
  Lemma~\ref{lem:signs}(ii)--(iv): near $0$ the last inequality of~(iii) gives $\re\tvphi>0$ on
  $\{\delta\le|z|\le\delta_0,\ \frac34\thlip\le|\arg z|\le\thw\}$ if $\delta_0$ is small, so the lips can descend
  to a small height there, follow the band at a height where~(ii) applies, and rise to $4c\pm ih_0$ along
  $|z-4c|=h_0$, where~(iv) applies). They stay outside $U(4c,h_0)$ except at their endpoint. We take $\delta$ so small that $\frac34\thlip\le|\arg z|\le\frac54\thlip$ on
  $\Sigma^\pm_{\mathrm{lip}}\cap U(0,\delta)$.
\item The \emph{vertical segments} $\Sigma_v^\pm:=\{4c\pm iy:\ 0<y\le h_0\}$.
\item The \emph{saturated lips} $\Sigma^{\triangle\pm}:=\{x\pm ih_0:\ 4c\le x\le n^p+\delta'\}$.
\end{itemize}
These arcs bound the \emph{lens} $\Ob^\pm$ (between the band and $\Sigma^\pm_{\mathrm{lip}}$, left of
$\Sigma^\pm_v$) and the \emph{saturated regions} $\Osat_\pm:=\{4c<\re z<n^p+\delta',\ 0<\pm\im z<h_0\}$. The rest of
$\C$ is $\Oinf$. We write $\Sigma$ for the union of all these arcs. Only finitely many arcs meet at each vertex, and
for fixed $n$, $\Sigma$ is bounded. In $U(4c,\eps)$ the only arcs are the band, the saturated segment and
$\Sigma^\pm_v$; at $0$ only the band and the two lips meet. Nothing lies on $(-\infty,0)$ or on
$(n^p+\delta',\infty)$.

The closures of $\Ob^\pm$ and $\Osat_\pm$, minus the point $0$, lie in $\Zn$. Indeed, near $0$ the lens lies in
$|\arg z|<\thw$ by construction of the lips; every point of $\Osat$ has $\re z\ge4c$ and $|\im z|\le h_0<4c\tan\thw$,
so $|\arg z|<\thw$; and $|z|\le n^p+\delta'+h_0
<e^{pc_0}n^p$ for $n\ge n_0$.

\begin{figure}[ht]
\centering
\begin{tikzpicture}[scale=0.95,>=Stealth]
  % panel 1: p >= 1
  \begin{scope}
    \draw[gray!60] (-1.6,0)--(0,0);
    \draw[thick] (0,0)--(1.8,0) node[below left]{\small band};
    \draw[thick] (0,0)--(60:1.5) node[above]{\small lip};
    \draw[thick] (0,0)--(-60:1.5);
    \draw[dashed] (0,0)--(120:1.3);
    \draw[dashed] (0,0)--(-120:1.3);
    \draw[gray] (0,0) circle (1.5);
    \node at (0,-2.0) {\small $p\ge1$: lips at $\pm\pi/3$};
    \node[gray] at (-1.0,0.25) {\tiny no jump};
  \end{scope}
  % panel 2: p < 1
  \begin{scope}[xshift=5cm]
    \draw[gray!60] (-1.6,0)--(0,0);
    \draw[thick] (0,0)--(1.8,0);
    \draw[thick] (0,0)--(22.5:1.5) node[right]{\small lip};
    \draw[thick] (0,0)--(-22.5:1.5);
    \draw[densely dotted] (0,0)--(45:1.6) node[above right]{\small seam};
    \draw[densely dotted] (0,0)--(-45:1.6);
    \draw[gray] (0,0) circle (1.5);
    \node at (0,-2.0) {\small $p<1$: lips at $\pm\pi p/4$, seams at $\pm\pi p/2$};
    \node[gray] at (-0.9,0.9) {\small $\Sout$};
    \node[gray] at (0.9,0.25) {\tiny $\Sin$};
  \end{scope}
  % panel 3: global
  \begin{scope}[xshift=9.4cm,yshift=0cm]
    \draw[thick] (0,0)--(4.4,0);
    \draw[thick] (0,0) .. controls (0.5,0.6) and (1.2,0.6) .. (1.6,0.6);
    \draw[thick] (0,0) .. controls (0.5,-0.6) and (1.2,-0.6) .. (1.6,-0.6);
    \draw[thick] (1.6,-0.6)--(1.6,0.6);
    \draw[thick] (1.6,0.6)--(4.0,0.6)--(4.0,-0.6)--(1.6,-0.6);
    \fill (0,0) circle (1pt) node[below left]{\small $0$};
    \fill (1.6,0) circle (1pt) node[below right]{\small $4c$};
    \fill (3.8,0) circle (1pt) node[below]{\small $n^p$};
    \node at (0.85,0.18) {\tiny $\Ob^+$};
    \node at (2.8,0.3) {\tiny $\Osat_+$};
    \node at (2.4,1.0) {\small $\Oinf$};
    \node at (2.2,-2.0) {\small all $p$ (not to scale)};
  \end{scope}
\end{tikzpicture}
\caption{The contour near the hard edge for $p\ge1$ and $p<1$ (Bessel rays dashed), and the whole contour. The
negative axis and, for $p<1$, the seams carry no jump of $V$ or $Q$. The seams are used only to construct the outer
scalar (\S\ref{sec:outer-plt1}).}
\label{fig:geometry}
\end{figure}

\subsection{The transformations}\label{sec:frame-transf}
We now remove the poles, normalise at infinity and open the lens. The pattern is that of \ref{tool:T1}(b); since we
work in a one-sided variable, we write every factor out. Put, in $\C_\pm$,
\begin{equation}\label{eq:kappa}
  \kappa_\pm:=\frac1{1-e^{\mp2i\theta}},\qquad\text{so that}\qquad \kappa_\pm-1=e^{\mp2i\theta}\kappa_\pm,
\end{equation}
and note $|e^{\mp2i\theta}|<1$ in $\C_\pm\cap\Zn$ by~\eqref{eq:imtheta}, so $\kappa_\pm$ is analytic there. Define,
on $\Zn\cap\C_\pm$,
\begin{equation}\label{eq:LA}
  L_\pm:=\begin{pmatrix}1&\mp We^{\mp2i\theta}\kappa_\pm\\0&1\end{pmatrix},\qquad
  A_\pm:=\begin{pmatrix}1&0\\\mp e^{2n\varphi}/W&1\end{pmatrix},\qquad
  A^\triangle_\pm:=\begin{pmatrix}1&0\\\mp1/(\kappa_\pm W)&1\end{pmatrix}.
\end{equation}

\emph{Step 1: removal of the poles.} Let
\begin{equation}\label{eq:R}
  R:=U\ \text{in }\Oinf,\qquad R:=UL_\pm\ \text{in }\Ob^\pm,\qquad R:=UA^\triangle_\pm\,\kappa_\pm^{\sigma_3}\ \text{in }\Osat_\pm .
\end{equation}
The restriction of $R$ to each of the four regions $\Ob^\pm$, $\Osat_\pm$ extends continuously to the real axis,
including the nodes. Indeed, near $X_k$ all factors are meromorphic, and $1-e^{\mp2i\theta}=\pm2i\theta'(X_k)
(z-X_k)(1+O(z-X_k))$ because $e^{2i\theta(X_k)}=1$. In $\Ob^\pm$ the entry $\mp We^{\mp2i\theta}\kappa_\pm$ then has the
simple pole $W(X_k)/(2i\theta'(X_k)(z-X_k))=w(X_k)/(z-X_k)$, and $R_{j2}=U_{j2}-U_{j1}w(X_k)/(z-X_k)+O(1)$ is regular
by Lemma~\ref{lem:interp}. In $\Osat_\pm$, $\kappa_\pm^{-1}$ vanishes simply at $X_k$, which removes the pole of the
second column; the first column is $\kappa_\pm(U_{j1}\mp U_{j2}/(\kappa_\pm W))$, and the bracket vanishes at $X_k$
since $U_{j2}/(\kappa_\pm W)\to\pm U_{j1}(X_k)$, so the first column is regular too.

\emph{Step 2: normalisation.} $T:=e^{-n\gamma_p\sigma_3}Re^{n\varphi\sigma_3}$.

\emph{Step 3: the lens.} $S:=TA_\pm$ in $\Ob^\pm$, and $S:=T$ elsewhere.

\emph{Step 4: the outer scalar.} Given a function $F$ with the properties (O1)--(O2) of~\S\ref{sec:frame-outer}
below and its lens version $\tilde F$, let
\begin{equation}\label{eq:V}
  V:=C_\infty\,S\,\tilde F^{\sigma_3}\ \text{in }\Ob\cup\Osat,\qquad V:=C_\infty\,S\,F^{\sigma_3}\ \text{in }\Oinf,
\end{equation}
with a constant diagonal matrix $C_\infty$ fixed in Lemma~\ref{lem:normalisation}.

Composing the four steps, and using $\kappa_\pm^{\sigma_3}\tilde F^{\sigma_3}=F^{\sigma_3}$ in $\Osat_\pm$
(definition~\eqref{eq:Ftilde} below),
\begin{equation}\label{eq:VM}
  V=C_\infty e^{-n\gamma_p\sigma_3}\,U\,M,\qquad
  M:=\begin{cases}e^{n\varphi\sigma_3}F^{\sigma_3}&\text{in }\Oinf,\\
  L_\pm e^{n\varphi\sigma_3}A_\pm\tilde F^{\sigma_3}&\text{in }\Ob^\pm,\\
  A^\triangle_\pm e^{n\varphi\sigma_3}F^{\sigma_3}&\text{in }\Osat_\pm.\end{cases}
\end{equation}
Since $U$ is analytic off the nodes, the jump of $V$ on any arc of $\Sigma$ away from the nodes is $M_-^{-1}M_+$.

The factor in $\Osat_\pm$ is the product of the factors $\Pi^{-\sigma_3}(\cdot)$ and $V^\triangle_\pm$ of \ref{tool:T1}(b),
moved into one step, except that in $\Osat_-$ we have changed it by the scalar $-1$. With the literal product a constant
jump $-I$ would appear on the saturated segment.

\subsection{Where the weight is evaluated}\label{sec:frame-where}
The following bookkeeping explains why the three strips $\vartheta_p$ of Table~\ref{tab:geometry} differ.

\begin{lemma}[Where the weight is evaluated]\label{lem:weightdomain}
The functions $w$, $W$, $r_n$ and $W_0$ (defined in~\eqref{eq:W0}) are evaluated in this paper only at points of
$\Zn$: in the closures of $\Ob$ and $\Osat$ (minus the point $0$), which contain the band, the saturated segment,
the lips, the vertical segments, $\Sigma^{\triangle\pm}$ and the top segment, and at the points of the lips and of the sectors $\{|\arg z|\le\frac54\thlip\}\cap U(0,\delta)$ used in
\S\ref{sec:hardedge} and Appendix~\ref{app:apriori}, and on the top segment. The region $\Oinf$ never uses the
weight: there $M=e^{n\varphi\sigma_3}F^{\sigma_3}$ contains neither $w$ nor $W$. Consequently, for every such point
$\xi(z)=\log t(z)$ lies in $\{\re\xi\ge-c_0,\ |\im\xi|\le\thw/p\}\subset S_{\vartheta_p}$.
\end{lemma}
% src: ROUTE2 §2 (H_an) note (REVIEW_Plan48 F4); ROUTE2_plt1 §1 Lemma S_w; ROUTE2_xframe1 §2 "Strip needed"

\begin{proof}
By~\eqref{eq:VM}, $w$ and $W$ enter $V$ only in $\Ob\cup\Osat$, whose closure minus $\{0\}$ lies in $\Zn$
(\S\ref{sec:frame-contour}); the top segment has $|z|\le n^p+\delta'+h_0<e^{pc_0}n^p$ and $|\arg z|<\thw$. The sectors
used near $0$ lie in $\Zn$ because $\frac54\thlip<\thw$. For $z\in\Zn$, $\re\xi=\log n-\frac1p\log|z|>-c_0$ and
$|\im\xi|=|\arg z|/p<\thw/p$.
\end{proof}

For $p>1$ every argument up to $\thw<\pi$ would be allowed, and the hypothesis is the full strip $\pi/p$. For $p<1$
the same bookkeeping keeps all evaluations in the inner sector $|\arg z|<\pi p/2$, so the strip $\pi/2$ suffices; the
strip $\pi/p$ would be wider than needed, and wider than the strip on which, for example,
$\Lambda(\xi)=\sqrt{\log(e^\xi+e)}$ is analytic. At $p=1$ the strip is $\pi/2$ for the same reason.

\begin{lemma}[Potter bounds in the sector]\label{lem:pottersector}
For every $\eta>0$ there is $C_\eta$ such that:
\begin{enumerate}
\item[(i)] $|\Lambda(\xi)|\le C_\eta+\eta\re\xi$ for $\xi\in S_{\vartheta_p}$, $\re\xi\ge0$, and $|\Lambda|\le C$ on
  $S_{\vartheta_p}\cap\{\re\xi\le0\}$. Hence $|\ell_1(t)|^{\pm1}\le C_\eta|t|^\eta$ for $|\arg t|\le\vartheta_p$,
  $|t|\ge1$.
\item[(ii)] $|r_n(z)|^{\pm1}\le C_\eta n^{2\eta}(1+|z|^{-\eta/p})$ for $z\in\Zn$.
\item[(iii)] For $R\ge1$, $z\in\Zn$ with $|z|\le R$ and $x\in(0,R]$,
  \begin{gather*}
    |\log r_n(z)-\log r_n(x)|\le\frac{\eps_\Lambda(T)}p\bigl(|\arg z|+|\log(x/|z|)|\bigr),\\
    |\log r_n(z)|\le\frac{\eps_\Lambda(T)}p\bigl(|\arg z|+|\log|z||\bigr),
  \end{gather*}
  with $T:=\log n-\frac1p\log R$. Here $\log r_n=\Lambda(\xi)-\Lambda(\log n)$.
\end{enumerate}
\end{lemma}
% src: ROUTE2 §2 "Consequences ... Potter in the sector"; ROUTE2 §3.3 (path); ROUTE2_plt1 §1

\begin{proof}
(i) For $\re\xi\ge T_0$, integrate $\Lambda'$ along the horizontal segment from $T_0+i\im\xi$ to $\xi$:
$|\Lambda(\xi)-\Lambda(T_0+i\im\xi)|\le\eps_\Lambda(T_0)\re\xi$. Choose $T_0$ with $\eps_\Lambda(T_0)\le\eta$; $\Lambda$ is
bounded on the compact set $S_{\vartheta_p}\cap\{0\le\re\xi\le T_0\}$ and, by hypothesis, on $\re\xi\le0$. Then
$|\ell_1(t)|^{\pm1}=e^{\pm\re\Lambda(\log t)}$.
(ii) $|r_n|^{\pm1}=|\ell_1(t)|^{\pm1}\ell_1(n)^{\mp1}$ with $|t|=n|z|^{-1/p}$; apply~(i) to both factors (for $|t|\le1$ use
the bound on $\re\xi\le0$).
(iii) Join $\xi(x)$ to $\xi(z)$ by the horizontal segment from $\log n-\frac1p\log x$ to $\log n-\frac1p\log|z|$
followed by the vertical segment to $\xi(z)$. Along this path $\re\xi\ge T$ and $|\im\xi|\le|\arg z|/p$, and its length
is $(|\log(x/|z|)|+|\arg z|)/p$. The second inequality is the case $x=1$, since $r_n(1)=1$.
\end{proof}

\subsection{The outer scalar and the jumps of $V$}\label{sec:frame-outer}
\S\ref{sec:outer} constructs, for each geometry, a function $F$ with the following five properties. The constants in
(O3) and (O4) do not depend on $n$; $\nu_p:=\min(1/p,\frac12)$, and $\eta_0\in(0,\delta]$ is fixed.
\begin{enumerate}
\item[(O1)] $F$ is analytic and zero-free in $\C\setminus[0,\infty)$ (in particular, it has no jump on $(-\infty,0)$),
  it has continuous boundary values on $(0,\infty)\setminus\{X_k\}$, and $F_+=F_-e^{2i\theta}$ on
  $(4c,\infty)\setminus\{X_k\}$.
\item[(O2)] The \emph{lens scalar}
  \begin{equation}\label{eq:Ftilde}
    \tilde F:=(1-e^{\mp2i\theta})F=F/\kappa_\pm\qquad\text{in }\C_\pm\cap\Zn
  \end{equation}
  extends to an analytic, zero-free function on $\Zn\setminus[4c,\infty)$, with $\tilde F_+=-\tilde F_-$ on
  $(4c,e^{pc_0}n^p)$. Thus $\tilde F^2$ is analytic and zero-free on $\Zn$, and so is the \emph{band factor}
  \begin{equation}\label{eq:Hn}
    H_n(z):=\frac1{(4c-z)\tilde F(z)^2}\qquad(z\in\Zn).
  \end{equation}
  We require $H_n>0$ on $(0,e^{pc_0}n^p)$ and that $H_n(0^+):=\lim_{z\to0,\,z\in\Zn}H_n(z)$ exists in $(0,\infty)$.
\item[(O3)] For $z\in\Zn$ with $|z|\le\eta_0$:
  $|\log H_n(z)-\log H_n(0^+)|\le C|z|^{\nu_p}$ and $|(\log H_n)'(z)|\le C|z|^{\nu_p-1}$.
\item[(O4)] For $z\in\Zn$ with $|z|\ge\eta_0$: $|\log H_n(z)-\log H_n(0^+)|\le C\bigl(1+\log(1+|z|)\bigr)^2$.
\item[(O5)] (For each fixed $n$.) $|F|^{\pm1}$ is bounded on $U(0,\delta)\setminus(\overline{\Ob}\cup[0,\infty))$, and
  $\tilde F^{\pm1}$ is bounded on $U(0,\delta)\cap\Ob$. There is $m_n$ with $|F(z)|^{\pm1}\le C_n|z|^{m_n}$ for
  $|z|\ge n^p+\delta'$, $z\notin\R$, and $F(iy)$ is bounded and bounded away from $0$ as $y\to\pm\infty$.
\end{enumerate}
The \emph{band weight} is
\begin{equation}\label{eq:W0}
  W_0:=\frac W{\tilde F^2}=C_n\,z^\alpha(4c-z)H_n(z)r_n(z)\quad\text{on }\Zn,\qquad C_n:=C_n'
\end{equation}
(Lemma~\ref{lem:bandweight}). On $(0,4c)$ it equals $C_n$ times a positive function.

\begin{lemma}[Jumps of $V$]\label{lem:jumpsV}
Assume (O1) and (O2). Then $V$ is analytic in $\C\setminus\Sigma$, has continuous boundary values on $\Sigma$ minus its
vertices (including at the nodes), and $V_+=V_-v_V$ with
\begin{enumerate}
\item[(a)] on the band $(0,4c)$: $v_V=\begin{pmatrix}0&W_0\\-1/W_0&0\end{pmatrix}$;
\item[(b)] on $\Sigma^\pm_{\mathrm{lip}}$ (oriented away from $0$), with $\epsilon_+=1$ and $\epsilon_-=0$:
  \[
    v_V=\begin{pmatrix}1&0\\e^{2n\varphi}/W_0&1\end{pmatrix}
    +\begin{pmatrix}\epsilon_\pm(\kappa_\pm-1)&W_0e^{-2n\tvphi}\\(\kappa_\pm-1)e^{2n\varphi}/W_0&\epsilon_\mp(\kappa_\pm-1)\end{pmatrix};
  \]
\item[(c)] on $\Sigma_v^\pm$ (oriented away from the real axis): $v_V=\begin{pmatrix}1&-W_0e^{-2n\tvphi}\\0&1\end{pmatrix}$;
\item[(d)] on the saturated segment $(4c,n^p+\delta')$: $v_V=\begin{pmatrix}1&0\\e^{2n\tvphi}/W_0&1\end{pmatrix}$;
\item[(e)] on $\Sigma^{\triangle\pm}$ (oriented from left to right): $v_V=\begin{pmatrix}1&0\\\kappa_\pm e^{2n\varphi}/W_0&1\end{pmatrix}$;
\item[(f)] on $\Sigma_{\mathrm{top}}\cap\C_\pm$ (oriented upwards): $v_V=\begin{pmatrix}1&0\\\mp\kappa_\pm e^{2n\varphi}/W_0&1\end{pmatrix}$;
\item[(g)] there is no jump on $(-\infty,0)$ or on $(n^p+\delta',\infty)$.
\end{enumerate}
\end{lemma}
% src: RH_px §3 "Facts about the jumps", §4.2, §4.3; RH_zeta_close §1.0; WLXZ (6.4) (T1); all recomputed here.

\begin{proof}
We use $M_-^{-1}M_+$ from~\eqref{eq:VM} and the identities
\begin{equation}\label{eq:identities}
  e^{\mp2i\theta}e^{-2n\varphi}=e^{-2n\tvphi},\qquad F=\kappa_\pm\tilde F,\qquad W/\tilde F^2=W_0\qquad\text{in }\C_\pm .
\end{equation}
For a diagonal matrix $D=\mathrm{diag}(d_1,d_2)$, conjugation $D^{-1}XD$ multiplies the $(1,2)$ entry by $d_2/d_1$ and the
$(2,1)$ entry by $d_1/d_2$.

(a) $L_-^{-1}L_+=\bigl(\begin{smallmatrix}1&W\\0&1\end{smallmatrix}\bigr)$: its $(1,2)$ entry is
$-W(e^{-2i\theta}\kappa_++e^{2i\theta}\kappa_-)=-W\bigl(\frac{e^{-2i\theta}}{1-e^{-2i\theta}}+\frac1{e^{-2i\theta}-1}\bigr)=W$. Conjugating by $e^{n\varphi_\pm\sigma_3}$ and using $\varphi_-=-\varphi_+$
(Lemma~\ref{lem:gfunction}(i)) gives $\bigl(\begin{smallmatrix}e^{2n\varphi_+}&W\\0&e^{2n\varphi_-}\end{smallmatrix}\bigr)$, and
\[
  A_-^{-1}\begin{pmatrix}e^{2n\varphi_+}&W\\0&e^{2n\varphi_-}\end{pmatrix}A_+=\begin{pmatrix}0&W\\-1/W&0\end{pmatrix}.
\]
Since $\tilde F$ has no jump on the band (O2), the final conjugation gives~(a) by~\eqref{eq:identities}.

(b) Upper lip: the $+$ side is $\Oinf$, so $v_V=M_{\Ob}^{-1}M_{\Oinf}=\tilde F^{-\sigma_3}A_+^{-1}\bigl(e^{-n\varphi\sigma_3}
L_+^{-1}e^{n\varphi\sigma_3}\bigr)F^{\sigma_3}$. The bracket is $\bigl(\begin{smallmatrix}1&W\kappa_+e^{-2n\tvphi}\\0&1
\end{smallmatrix}\bigr)$, and
\[
  A_+^{-1}\begin{pmatrix}1&W\kappa_+e^{-2n\tvphi}\\0&1\end{pmatrix}=
  \begin{pmatrix}1&W\kappa_+e^{-2n\tvphi}\\e^{2n\varphi}/W&\kappa_+\end{pmatrix},
\]
because $1+\kappa_+e^{2n\varphi}e^{-2n\tvphi}=1+\kappa_+e^{-2i\theta}=\kappa_+$. Conjugating with $\tilde F^{-\sigma_3}(\cdot)F^{\sigma_3}$,
$F=\kappa_+\tilde F$, gives $\bigl(\begin{smallmatrix}\kappa_+&W_0e^{-2n\tvphi}\\\kappa_+e^{2n\varphi}/W_0&1\end{smallmatrix}\bigr)$,
which is~(b). The lower lip (where the $+$ side is the lens) is the mirror computation and gives
$\bigl(\begin{smallmatrix}1&W_0e^{-2n\tvphi}\\\kappa_-e^{2n\varphi}/W_0&\kappa_-\end{smallmatrix}\bigr)$.

(c) Upper segment, $+$ side $=\Ob^+$: $v_V=M_{\Osat}^{-1}M_{\Ob}=F^{-\sigma_3}e^{-n\varphi\sigma_3}(A^\triangle_+)^{-1}L_+
e^{n\varphi\sigma_3}A_+\tilde F^{\sigma_3}$. One checks $(A^\triangle_+)^{-1}L_+=\bigl(\begin{smallmatrix}1&-We^{-2i\theta}
\kappa_+\\1/(\kappa_+W)&1-e^{-2i\theta}\end{smallmatrix}\bigr)$, and after the conjugation by $e^{n\varphi\sigma_3}$ and right
multiplication by $A_+$ the $(2,1)$ entry vanishes and the diagonal becomes $(\kappa_+,\kappa_+^{-1})$; the final
conjugation with $F=\kappa_+\tilde F$ removes the diagonal and gives~(c). The lower segment is the mirror image.

(d) On $(4c,n^p+\delta')$ both sides lie in $\Osat$. Since $\kappa_\pm^{-1}=1-e^{\mp2i\theta}$,
$(A^\triangle_-)^{-1}A^\triangle_+=\bigl(\begin{smallmatrix}1&0\\-(\kappa_+^{-1}+\kappa_-^{-1})/W&1\end{smallmatrix}\bigr)$, and
$\kappa_+^{-1}+\kappa_-^{-1}=2-2\cos2\theta=\kappa_+^{-1}\kappa_-^{-1}$. Conjugating by $e^{n\varphi\sigma_3}F^{\sigma_3}$ with
boundary values, the diagonal is $(e^{n(\varphi_+-\varphi_-)}F_+/F_-)^{\sigma_3}=(e^{-2i\theta}e^{2i\theta})^{\sigma_3}=I$
by Lemma~\ref{lem:gfunction}(i) and (O1), and the $(2,1)$ entry is
$-e^{n(\varphi_++\varphi_-)}F_+F_-/(\kappa_+\kappa_-W)=-e^{2n\tvphi}\tilde F_+\tilde F_-/W$. By (O2),
$-\tilde F_+\tilde F_-=\tilde F^2$, which gives~(d).

(e), (f) Here $v_V=(F^{-\sigma_3}e^{-n\varphi\sigma_3})(A^\triangle_\pm)^{\pm1}(e^{n\varphi\sigma_3}F^{\sigma_3})$ with the
exponent fixed by the orientation, and the $(2,1)$ entry $\mp1/(\kappa_\pm W)$ of $A^\triangle_\pm$ becomes
$\mp F^2e^{2n\varphi}/(\kappa_\pm W)=\mp\kappa_\pm e^{2n\varphi}/W_0$.

(g) On $(-\infty,0)$: $U$ has no jump, $e^{n\varphi}$ jumps by $e^{-2\pi in}=1$ (Lemma~\ref{lem:gfunction}(i)) and $F$ has no
jump (O1). On $(n^p+\delta',\infty)$ there are no nodes, and $M_-^{-1}M_+=(e^{n(\varphi_+-\varphi_-)}F_+/F_-)^{\sigma_3}=I$ as
in~(d).

Continuity at the nodes: in~(a) and~(d) all entries are analytic across the real axis; together with Step~1 this shows
that $V$ has continuous boundary values there.
\end{proof}

\begin{table}[ht]
\centering\small
\renewcommand{\arraystretch}{1.25}
\begin{tabular}{@{}p{3.4cm}p{3.2cm}p{3.2cm}p{3.6cm}@{}}
\toprule
& $p>1$ & $p=1$ & $0<p<1$\tabularnewline
\midrule
\raggedright $\arg t$ on $\C\setminus(-\infty,0]$ & \raggedright $(-\pi/p,\pi/p)\subset(-\pi,\pi)$ & \raggedright $(-\pi,\pi)$, open & \raggedright wider than $(-\pi,\pi)$; poles of $\Gamma(t)$ on $\arg z=\pm\pi p$\tabularnewline
\raggedright where $\tilde D=e^{-j(t)}$ is used & \raggedright slit plane & \raggedright $|\arg z|\le\thw=7\pi/16$ & \raggedright inner sector $\Sin=\{|\arg z|<\pi p/2\}$\tabularnewline
\raggedright outer scalar $F$ & \raggedright $DBK$, Binet-cut scalar $B$ (Lemma~\ref{lem:binetcut}) & \raggedright $DK$, $D=e^{j(-n/z)}$, no cut (Lemma~\ref{lem:reflection}) & \raggedright $DB_{\mathrm{in}}Ke^{\chi}$ in $\Sin$, $B_{\mathrm{out}}Ke^{\chi}$ in $\Sout$ (Lemma~\ref{lem:seam})\tabularnewline
\raggedright $|\log\frac{H_n(z)}{H_n(0^+)}|$ near $0$ & \raggedright $\le C|z|^{1/p}$ & \raggedright $\le C|z|$ & \raggedright $\le C|z|(1+\log\frac1{|z|})$\tabularnewline
\raggedright $g(z)-\log z$ at $\infty$ & \raggedright $O(|z|^{-1/p})$ & \raggedright $O(|z|^{-1}\log|z|)$ & \raggedright $O(|z|^{-1})$\tabularnewline
\raggedright lip angle $\thlip$; Bessel rays & \raggedright $\pi/3$; $\pm2\pi/3$ & \raggedright $\pi/3$; $\pm2\pi/3$ & \raggedright $\pi p/4$; $\pm(\pi-\pi p/4)$\tabularnewline
\raggedright angle $\thw$ (weight evaluated) & \raggedright $7\pi/16$ & \raggedright $7\pi/16$ & \raggedright $3\pi p/8$\tabularnewline
\raggedright super-small entries $|e^{\mp2i\theta}|$ on the lips & \raggedright $e^{-2\pi n\sin(\frac{\pi}{3p})|z|^{-1/p}}$ & \raggedright $e^{-\sqrt3\pi n/|z|}$ & \raggedright $e^{-\sqrt2\pi n|z|^{-1/p}}$\tabularnewline
\raggedright strip $\vartheta_p$ for $\ell_1$ & \raggedright $\pi/p$ & \raggedright $\pi/2$ & \raggedright $\pi/2$\tabularnewline
\raggedright $\nu_p$ (in $\eps'_n=n^{-2\nu_p}\log n+\eps_\Lambda(T_n)$) & \raggedright $\min(1/p,\frac12)$ & \raggedright $\frac12$ & \raggedright $\frac12$\tabularnewline
\raggedright classical quotient parametrix at $0$ & \raggedright works for pure powers (not used) & \raggedright fails (Remark~\ref{rem:quotient-p1}) & \raggedright works with a seam glue (not used)\tabularnewline
\bottomrule
\end{tabular}
\caption{The ingredients of the Riemann--Hilbert analysis that depend on the geometry. Everything else in
\S\S\ref{sec:frame}--\ref{sec:extraction} is the same for all $p>0$. The super-small entries are evaluated on the rays
$\arg z=\pm\thlip$; $\eps'_n$ is the lip-mismatch size of \S\ref{sec:hardedge}.}
\label{tab:geometry}
\end{table}
